\caption{\textbf{Uncertain processes, selective evidence, and two sides of burden.}
For one jurisdiction--program--version, $N$ is the finite number of steps,
$G$ their identities and routing, and $\mathcal A$ the agency entities and aliases.
Step $i$ has actor $a_i$, burden parameters $\theta_i$, disclosure strength
$\lambda_i$, and visibility $v_i$. Shared agency effects $\alpha_g$ and task-type
effects $\beta_t$, governed by hyperparameters $\eta$, pool information across
activities and programs. Actor $a_i$ selects the applicable agency effect.
Case characteristics $X_n$ select a route and outcome $R_n$; each step visit
$E_{nir}$ realizes elapsed days, applicant work-hours, agency work-hours, fees and
encounters. Payer burden $B_n^u$ and agency burden $B_n^a$ are deterministic,
vector-valued summaries of those visits and their schedule. Document metadata
$D_m$ and reporting parameters $\psi$ govern which steps or segments $Z_m$ are
disclosed and what typed assertions $C_m$ the source makes. The $E\to Z_m$ edge
allows long or burdensome cases to be reported disproportionately. Extractor error
$\varepsilon$ connects those latent assertions to the observed package $Y_m$.
An official median is informed directly by $\theta_i$; a case account by
$E_{nir}$. Elapsed time is never equated with hands-on labor. The figure specifies
the target model; the accompanying glossary distinguishes implemented pilot
components, fixed assumptions, and unimplemented extensions.}
\label{fig:plate-simplified}
